% !TeX root = ../main.tex
% Add the above to each chapter to make compiling the PDF easier in some editors.

\chapter{Implementation}\label{chapter:implementation}

\section{Project Setup}

The projects as well as this thesis' source code are avaibable in the following Github repository:

\url{https://github.com/mulsicpp/bachelor_thesis_2025.git}

The project is organized in the following directories:

\begin{itemize}
    \item "assets": This folder contains the projects' assets. It is further divided into scenes, which contains the glTF test scenes, and "shaders", which contains the shaders, that are used in the different pipelines.
    \item "external": This folder contains all of the "complex""" external libraries. "Complex" in this context means, that the library has multiple header or source files or it is in binary form.
    \item "premake": This folder contains the premake executables for both windows and linux.
    \item "src": This folder contains all of our source files for the project. It is described in more detail in the next chapter.
\end{itemize}

Additionally, there are some scripts, that are used to build the project.

\subsection{Source Code Organization}

In this section we will specifically look at how the "src" folder is organised.

Directly in the "src" folder we find "main.cpp", which contains our entry point of the application.
The "external" folder once again contains external libraries, but this time the libraries, which are in the form of a single header file. These usually need to be included in a C/C++ file and need to be compiled. That's why it makes sense to keep these libraries with our other source files.
The "utils" folder contains utilities, that may be useful all throughout the project. It includes functionalities like debug logging, move semantics and file loading.
All the folders with a "vk\_" prefix contain code, which abstracts the low-level Vulkan API code. The goal of these abstractions is, that the low-level code should only ever be used in these folders. The rest of the project should only use the abstracted API. Each of the folders focuses on a different area of Vulkan:

\begin{itemize}
    \item "vk\_core": The most important class in this folder is the \texttt{Context} class. It is designed as a singleton, which contains the \texttt{VkInstance}, \texttt{VkPhysicalDevice}, \texttt{VkDevice}, \texttt{VkQueues} and \texttt{VmaAllocator}. After the context gets created it can be accessed anywhere. The folder also contains the \texttt{Handle<T>} class, which ensures that a handle only has one owner, the \texttt{CommandBuffer} abstraction and functionality to deals with \texttt{VkFormat}.
    \item "vk\_resources": This folder contains Vulkans resource objects like \texttt{Buffer} and \texttt{Image}, as well as \texttt{ImageView} and \texttt{SubBuffer}.
    \item "vk\_pipeline": This folder contains Vulkan objects, that are related to the setup of a pipeline. This includes \texttt{Pipeline}, \texttt{PipelineLayout}, \texttt{Shader}, \texttt{RenderPass}, \texttt{Framebuffer} and \texttt{DescriptorPool}.
    \item "vk\_sync": Contains the vulkan synchronization objects \texttt{Fence} and \texttt{Semaphore}.
    \item "scene": This folder contains the \texttt{Scene} class, which describes a dynamic scene with a scene graph. All the other classes used by or contained in the scene, like \texttt{Mesh}, \texttt{Node}, \texttt{Camera} and \texttt{Animation}, can also be found in this folder.
    \item "rendering": This folder contains 3 different renderers:
    \begin{itemize}
        \item \texttt{Rasterizer}: Render a scene from the viewpoint of a camera into a framebuffer using traditional rasterisation.
        \item \texttt{Raytracer}: Render a scene from the viewpoint of a camera into an image using modern hardware-accelerated raytracing.
        \item \texttt{Skinner}: Calculates the skinned positions of a mesh during an animation.
    \end{itemize}
\end{itemize}


\section{Move Semantics}
As already mentioned in the Methods and Technology chapter, Rust was originally considered as programming language, because it is modern and memory safe. One of Rust's core features is ownership, which means that each value/object can only be owned and managed by only one variable, and when that variable goes out of scope, it alone is responsible for cleaning up the object. When that variable gets assigned to another variable, the new variable becomes the new owner (i.e. the value gets moved). Ownership would be a very nice feature to manage \getVkObjects{}. \getVkObjects{} should not be copied on assignment but only moved and when the object is no longer needed (i.e. the owner goes out of scope), it should automatically be destroyed.
While C++ does not support ownership directly, it does have so-called move semantics, with which we can emulate ownership. The ownership of Vulkan handles is managed by the \texttt{Handle<T>} class. To do this the class declares a default constructor, move constructor and destructor and overrides the move assignment operator. It is also important, that the copy constructor and copy assignment operator are deleted. The reason for that is that copying a \getVkObject{} is usually very expensive or not even possible and it is almost never necessary. When we want to copy an object, we need to explicitly create a new one first. This avoids unnecessary and complicated accidental creation and copying of new objects.

\section{Shared Pointers}
Since C++11 the language supports so-called smart pointers, which make memory management safer and less complicated. One of these smart pointers is \texttt{std::shared\_ptr<T>} (or our alias \texttt{ptr::Shared<T>}). This type of pointer allocates an object, which can then be shared between multiple owners. The object lives while it has at least one owner, and the last object to go out of scope cleans up the object. This structure can be combined very nicely with our previously defined handle. A \getVkObject{} can only have one owner, which is a harsh limitation and often causes problems. If we wrap the object in a shared pointer, it can be shared across multiple owners. We also added the class \texttt{ToShared<T>}, which a class can inherit from. An instance of that class can then be moved into a shared pointer by calling \texttt{to\_shared()}.

\section{The Builder Pattern}
As mentioned, before we want the creation of \getVkObjects{} to be very explicit. The first method to create objects in C++ you think of are constructors, but they prove to be problematic for a couple of reasons. The first problem is that constructors are often called implicitly (e.g. the default constructor gets called when declaring a variable). Default construction only makes sense for very few \getVkObjects{} (like for example a VkFence) and we do not want a \getVkObject{} to be created when implicitly (or explicitly) calling the default constructor. So, for that reason, the default constructor simply does not create a \getVkObject{} or, in other words, leaves the handle at \texttt{VK\_NULL\_HANDLE}. Now if a different constructor actually creates an object, the behavior of constructors becomes inconsistent regarding \getVkObject{} creation. In addition to that, \getVkObjects{} are almost always created using a \texttt{*CreateInfo} struct, which may take a lot of parameters. To solve all of these issues we implemented a \texttt{Builder} class for most \getVkObjects{}. Some exceptions are \texttt{Fence} and \texttt{ImageView}, which are created with \texttt{Fence::create(...)} and \texttt{ImageView::create\_from(...)} respectively. This makes it very clear when a \getVkObject{} is created, since it can only be created using a \texttt{Builder} or a static method (like \texttt{Fence::create(...)}) and especially it can never be created by a constructor. This is also the reason \getVkObjects{} never have constructor, besides the default constructor.

\begin{lstlisting}[language=C++]

vk::Buffer buffer; // declare a buffer (handle is VK_NULL_HANDLE)

buffer = vk::BufferBuilder()
    .add_usage(VK_BUFFER_USAGE_TRANSFER_DST_BIT)
    .add_usage(VK_BUFFER_USAGE_VERTEX_BUFFER_BIT)
    .memory_usage(VMA_MEMORY_USAGE_GPU_ONLY)
    .size(128)
    .build();


\end{lstlisting}

\section{The Vulkan context}

Before we can actually do any work in Vulkan, we need to do a lot of setup. This usually includes the following steps:

\begin{enumerate}
    \item Create a \texttt{VkInstance}.
    \item Select a \texttt{VkPhysicalDevice}, which supports the extenstions, features and capabilities required by the project.
    \item Create a \texttt{VkDevice} from the selected \texttt{VkPhysicalDevice}.
    \item When creating the \texttt{VkDevice}, \texttt{VkQueue}s from a certain queue family need to be enabled, to be able to submit command buffers later.
\end{enumerate}

These \getVkObjects{} are accessed very frequently throughout the project. For example, you need to pass the \texttt{VkDevice} as a parameter when creating or destroying most \getVkObjects{} and a command buffer, which executes a sequence of Vulkan commands, needs to be submitted to a certain \texttt{VkQueue}. Keeping track of all of these \getVkObjects is quite tideous, which is why we introduced the Vulkan context. It is represented by the \texttt{vk::Context} class, which stores these \getVkObjects, allows the programmer to access them anywhere and reduces the effort of stting up Vulkan.

The class is designed similar to a singleton, but it deviates in some parts. A singleton usually creates its instance at the start of the application or when it is used for the first time (lazy instantiation). However, \texttt{vk::Context} can not be implicitly instanciated this way, because it needs some parameters stored in a \texttt{vk::ContextInfo}, when being instanciated. For this reason we need to explicitly create the Vulkan context using \texttt{vk::Context::create(...)} with the desired \texttt{vk::ContextInfo} before calling \texttt{vk::Context::instance()} for the first time or we will get a \texttt{nullptr}. While all other objects in our project get cleaned up automatically when going out of scope, the Vulkan context also needs to be destroyed explicitly using \texttt{vk::Context::destroy()} when it is no longer needed. If we now want to get the current \texttt{VkDevice}, we can call \texttt{vk::Context::instance()->get\_device()}.
